\subsection{Geometry of Metric Measure Spaces}
In this section, we present the framework of metric measure space geometry introduced in~\cite{sturm2006geometry}. We first recall that 
%the specific context of a 
connected compact Riemannian manifolds fall under this framework, and then we introduce Gromov-Wasserstein metrics between metric measure spaces.

\subsubsection{Metric Measure Spaces}

A {\it metric measure space} (or mm-space for short) is a triple $(\spa,\dis,\mes)$ where $(\spa,\dis)$ is a Polish metric space (i.e., complete and separable)
%, meaning that it is complete 
    %(i.e. Cauchy sequences converge) 
%    and separable
    % (i.e. contains a countable dense subset),
and where $\mes$ is a measure on the Borel $\sigma$-algebra of $(\spa,\dis)$ 
    %(i.e. the $\sigma-$algebra generated by the open sets induced by the metric $\dis$),
    that is locally finite.\footnote{By locally finite, we mean that $\forall x\in \spa,\,\exists r>0$ such that $\mes(\ball(x,r))<\infty$.}
%We will call a mm-space $(\spa,\dis,\mes)$ normalized if $\mes(\spa)=1$.
The support of $\mes$
%, i.e. $\{x\in \spa\,:\,\text{every open set containing }x\text{ has a positive measure}\}$, 
is denoted as $\supp(\mes)$.

\begin{proposition}
    A connected compact Riemannian manifold $(\man,\g)$ together with its geodesic distance $\dis$ and the volume measure $(\man,\dis,\vol)$ is a mm-space. 
\end{proposition}

\begin{proof}
Since $\man$ is connected and compact, %, the Hopf-Rinow theorem~(see for example Section VII, Theorem 7.7. in~\cite{boothby1986introduction}) proves that $(\man,\dis)$ 
it is a complete metric space. It is also separable since %, by definition, any topological manifold 
it is second-countable. $\man$ being compact, the volume measure
$\vol$ is locally finite.
\end{proof}

%We now introduce some useful notions concerning metric measure spaces, beginning with isomorphy.\\
Let $(\spa_1,\dis_1,\mes_1)$ and $(\spa_2,\dis_2,\mes_2)$ be two mm-spaces. Recall that a map $\Phi\colon\supp(\mes_1)\rightarrow\supp(\mes_2)$ is an isomorphism of mm-spaces between $(\spa_1,\dis_1,\mes_1)$ and $(\spa_2,\dis_2,\mes_2)$ if $\Phi$ is an isometry of metric spaces
%, i.e. a bijection that verifies $\forall x,y\in~\supp(\mes_1):$ $\dis_2(\Phi(x),\Phi(y))=\dis_1(x,y)$,
%   \item $\mes_2=\mes_1\circ\Phi^{-1}$, i.e.
    and if $\mes_2$ is the pushforward measure of $\mes_1$ under the map $\Phi$.
As isomorphisms induce an equivalence relation of mm-spaces, we denote the isomorphism equivalence class of a mm-space $(\spa,\dis,\mes)$ as $[\spa,\dis,\mes]$.




Metric measure space isomorphisms also induce the notion of {\it variance} of a mm-space $(\spa,\dis,\mes)$, defined as:
$$\var(\spa,\dis,\mes)=\inf_{\substack{(\spa',\dis',\mes')\in [\spa,\dis,\mes]\\ z\in \spa'}}\int_{\spa'}\dis'(x,z)^2\,d\mes'(x).$$
If $(\spa,\dis,\mes)$ is a compact mm-space then clearly one has $\inf_{z\in \spa}\int_{\spa}\dis(x,z)^2\,d\mes(x) < \infty $ and then $\var(\spa,\dis,\mes)$ is finite. Moreover, the variance is by definition an invariant of mm-space isomorphisms.

 


%\begin{proposition}
%    If $(\spa,\dis,\mes)$ is a normalized mm-space and is bounded as a metric space then $\var(\spa,\dis,\mes)$ is finite.
%\end{proposition}
%\begin{proof}
%Since $(\spa,\dis,\mes)\in[\spa,\dis,\mes]$,
%$$\var(\spa,\dis,\mes)\leq\inf_{z\in \spa}\int_{\spa}\dis(x,z)^2\,d\mes(x).$$
%However, $(\spa,\dis)$ is bounded and therefore there exists $C>0$ that satisfies $\forall x,z\in \spa$: $\dis(x,z)\leq C$. We conclude that:
%\begin{align*}
%    \var(\spa,\dis,\mes)&\leq C^2\cdot\mes(X)\\
%    &\leq C^2.
%\end{align*}
%\end{proof}


\subsubsection{Gromov-Wasserstein Distance}

\begin{definition}{Wasserstein distance.}\\
Let $(\spa,\dis)$ be a complete and separable metric space. Let $\mes_1$ and $\mes_2$ be two measures on the Borel $\sigma$-algebra of $(\spa,\dis)$. For $p \geq 1$, the {\it $p$-Wasserstein distance} between $\mes_1$ and $\mes_2$ is defined as:
$$\wasser_p(\mes_1,\mes_2)=\left(\inf_{\mes}\int_{\spa\times\spa} \dis(x,y)^p\,d\mes(x,y)\right)^{\frac{1}{p}},$$
where the infimum is taken over all measures $\mes$ on $\spa\times \spa$ that have marginals $\mes_1$ and $\mes_2$.
\end{definition}


\begin{definition}{Metric coupling.}\\
Let $(\spa_1,\dis_1,\mes_1)$ and $(\spa_2,\dis_2,\mes_2)$ be two mm-spaces.
A {\it metric coupling} of $\dis_1$ and $\dis_2$ is any pseudometric\footnote{A pseudometric can be zero for non-equal points.} $\dis$ on $\spa_1\sqcup \spa_2$ that satisfies:
$$\forall x,y\in \supp(\mes_1):\,\dis(x,y)=\dis_1(x,y),$$
$$\forall x,y\in \supp(\mes_2):\,\dis(x,y)=\dis_2(x,y).$$
\end{definition}

\begin{definition}{Gromov-Wasserstein distance.}\\
 For $p \geq 1$, the  {\it $p$-Gromov-Wasserstein distance} between two mm-spaces $(\spa_1,\dis_1,\mes_1)$ and $(\spa_2,\dis_2,\mes_2)$ is defined as:
$$\gromov_p((\spa_1,\dis_1,\mes_1),(\spa_2,\dis_2,\mes_2))=\left(\inf_{\dis,\mes}\int_{\spa_1\times \spa_2} \dis(x,y)^p\,d\mes(x,y)\right)^{\frac{1}{p}},$$
where the infimum is taken over all measures $\mes$ on $\spa_1\times \spa_2$ that have marginals $\mes_1$ and $\mes_2$, and over all metric couplings $\dis$ of $\dis_1$ and $\dis_2$.
\end{definition}

    For the sake of simplicity, we will denote $\gromov_p((\spa_1,\dis_1,\mes_1),(\spa_2,\dis_2,\mes_2))$ simply as $\gromov_p(\spa_1,\spa_2)$ when there is no ambiguity over the metric and the measure associated to each space.

\begin{proposition}\label{prop:gw_def}
The $p$-Gromov-Wasserstein distance between two mm-spaces, $(\spa_1,\dis_1,\mes_1)$ and $(\spa_2,\dis_2,\mes_2)$, can be equivalently defined as:
$$\gromov_p((\spa_1,\dis_1,\mes_1),(\spa_2,\dis_2,\mes_2))=\inf_{\varphi_1,\varphi_2}\wasser_p(\mes_1\circ\varphi_1^{-1},\mes_2\circ\varphi_2^{-1}),$$
where the infimum is taken over all isometric embeddings $\varphi_1\colon\supp(\mes_1)\rightarrow \spa$ and $\varphi_2\colon\supp(\mes_2)\rightarrow \spa$ of the supports of $\mes_1$ and $\mes_2$ into a common metric space $(\spa,\dis)$.
\end{proposition}

In~\cite{memoli2011gromov}, an alternative definition of the $p$-Gromov-Wasserstein distance, inspired by a reformulation of the Gromov-Hausdorff distance, is provided as:
$$\Hat{\gromov_p}((\spa_1,\dis_1,\mes_1),(\spa_2,\dis_2,\mes_2))=\frac 12 \left(\inf_{\mes}\int_{(\spa_1\times \spa_2)^2} |\dis_1(x,x')-\dis_2(y,y')|^p\,d\mes(x,y)\otimes\mes(x',y')\right)^{\frac{1}{p}},$$
where the infimum is taken over all measures $\mes$ on $\spa_1\times \spa_2$ that have marginals $\mes_1$ and $\mes_2$. Due to its computational advantage for practical applications (such as, e.g., optimization), this formulation is popular in the literature, see for example~\cite{zhang2023duality,memoli2009spectral}.
Moreover, it is shown in~\cite{memoli2011gromov}  that $\Hat{\gromov_p}\leq\gromov_p$ (see Theorem 5.1(g) therein). In this article, we provide upper bounds for $\gromov_p$ which therefore automatically transfer to $\Hat{\gromov_p}$.

\begin{proposition}\label{prop:gw_met}
The $2$-Gromov-Wasserstein distance $\gromov_2$ is a metric on the set of mm-space isomorphism classes of finite variance.    
\end{proposition}

For proofs of Propositions \ref{prop:gw_def} and \ref{prop:gw_met}, see \cite{sturm2006geometry}.